\newcommand{\dia}{\textsf{diam}}
\newcommand{\girth}{g}
\newcommand{\con}{\textnormal{\textsf{con}}}
\newcommand{\lcon}{\ell^{\con}}
\newcommand{\Vcon}{V^{\con}}
\subsection{Completing the Proof of \Cref{thm: 0-ext with Steiner lower}}
\label{subsec: mapping}

We have shown in \Cref{subsec: canonical} all canonical solutions with size $O(k\log^{1-\eps} n)$ have cost $\Omega(\eps n\log\log n)$. In this subsection, we complete the proof of \Cref{thm: 0-ext with Steiner lower} by showing that, intuitively, an arbitrary solution $(\fset,\delta)$ to the instance $(G,T,D)$ can be ``embedded'' into a canonical solution, without increasing its cost by too much. 

We start by introducing the notion of \emph{continuization}.

\paragraph{Continuization of a graph.}
Let $G=(V,E,\ell)$ be an edge-weighted graph. Its \emph{continuization} is a metric space $(V^{\con},\ell^{\con})$, that is defined as follows. 
Each edge $(u,v)\in E$ is viewed as a continuous line segment $\con(u,v)$ of length $\ell_{(u,v)}$ connecting $u,v$, and the point set $V^{\con}$ is the union of the points on all lines $\set{\con(u,v)}_{(u,v)\in E}$. 
Specifically, for each edge $(u,v)\in E$, the line $\con(u,v)$ is defined as \[\con(u,v)=\set{(u,\alpha)\mid 0\le \alpha\le \ell_{(u,v)}}=\set{(v,\beta)\mid 0\le \beta\le \ell_{(u,v)}},\]
where $(u,\alpha)$ refers to the unique point on the line that is at distance $\alpha$ from $u$, and $(v,\beta)$ refers to the unique point on the line that is at distance $\beta$ from $v$, so $(u,\alpha)=(v,\ell_{(u,v)}-\alpha)$.

The metric $\ell^{\con}$ on $V^{\con}$ is naturally induced by the shortest-path distance metric $\dist_{\ell}(\cdot,\cdot)$ on $V$ as follows.
For a pair $p,p'$ of points in $V^{\con}$, 
\begin{itemize}
\item if $p,p'$ lie on the same line $(u,v)$, say $p=(u,\alpha)$ and $p'=(u,\alpha')$, then $\lcon(p,p')=|\alpha-\alpha'|$;
\item if $p$ lies on the line $(u,v)$ with $p=(u,\alpha)$ and $p'$ lies on the line $(u',v')$ with $p'=(u',\alpha')$, then 
\[
\begin{split}
\lcon(p,p')=\min\{
&  \dist_{\ell}(u,u')+\alpha+\alpha', \quad
	\dist_{\ell}(u,v')+\alpha+(\ell_{(u',v')}-\alpha'),\\
&	\dist_{\ell}(v,u')+(\ell_{(u,v)}-\alpha)+\alpha', \quad
	\dist_{\ell}(v,v')+(\ell_{(u,v)}-\alpha)+(\ell_{(u',v')}-\alpha')
\}.
\end{split}
\]
\end{itemize}
Clearly, every vertex $u\in V$ also belongs to $\Vcon$, and for every pair $u,u'\in V$, $\dist_{\ell}(u,u')=\lcon(u,u')$.
For a path $P$ in $G$ connecting $u$ to $u'$, it naturally corresponds to a set $P^{\con}$ of points in $\Vcon$, which is the union of all lines corresponding to edges in $E(P)$. The set $P^{\con}$ naturally inherits the metric $\lcon$ restricted on $P^{\con}$. We will also call $P^{\con}$ a path in the continuization $(\Vcon,\lcon)$.

We show that, for each graph $G$ with a set $T$ of terminals, then any other metric $w$ on a set of points containing $T$, such that $w$ restricted on $T$ is identical to $\dist_{G}$ restricted on $T$, can be ``embedded'' into the continuation of $G$, with expected stretch bounded by some structural measure that only depends on $G$. 
Specifically, we prove the following main technical lemma.

%In this section, we give a way to embed any metric into a graph $G$.

\begin{lemma} \label{lem:d-g}
Let $(G,T,\ell)$ be any instance of $\zesn$ such that $G$ is not a tree (so $\gir(G)<+\infty$), and let $(\fset,\delta)$ be any solution to it. Let $(\Vcon,\lcon)$ be the continuization of graph $G$.
Then there exists a random mapping $\phi: \fset \to \Vcon$, such that
\begin{itemize}
\item for each terminal $t\in T$, if $F$ is the (unique) cluster in $\fset$ that contains $t$, then $\phi(F)=t$; and
\item for every pair $F,F'\in \fset$,
%we can randomly construct a projection $f$ from the nodes in $\delta$ to the veritces in $V$ such that: for any pair of nodes $F_1$ and $F_2$,
$$\ex{\lcon(\phi(F),\phi(F'))} \le O\bigg(\frac{\diam(G)}{\gir(G)}\bigg)\cdot  \delta(F,F').$$
\end{itemize}
\end{lemma}

Before we prove \Cref{lem:d-g} in \Cref{subsec: diam-girth}, we provide the proof of \Cref{thm: 0-ext with Steiner lower} using it.

\emph{Proof of \Cref{thm: 0-ext with Steiner lower}.}
Consider any solution $(\fset,\delta)$ to the instance $(G,T,\ell)$ constructed in \Cref{subsec: instance} with size $|\fset|\le o(k\cdot \log^{1-\eps}k)$. 
From \Cref{obs: girth} and \Cref{cor: diam}, $\frac{\diam(G)}{\gir(G)}\le \frac{O(\log n)}{(\log n)/100}=O(1)$. From \Cref{lem:d-g}, there exists a random mapping $\phi: \fset \to \Vcon$, such that for every pair $F,F'\in \fset$,
%we can randomly construct a projection $f$ from the nodes in $\delta$ to the veritces in $V$ such that: for any pair of nodes $F_1$ and $F_2$,
$\ex{\lcon(\phi(F),\phi(F'))} \le O(1)\cdot  \delta(F,F')$. 

Fix such a mapping $\phi$, we define a canonical solution $(\fset,\hat \delta)$ based on $(\fset,\delta)$ as follows. The collection of clusters is identical to the collection $\fset$. For each $F\in \fset$, recall that $\phi(F)$ is a point in $\Vcon$. Assume the point $\phi(F)$ lies on the line $(u,v)$ and is closer to $u$ than to $v$ (i.e., $\lcon(\phi(F),u)\le\lcon(\phi(F),v)$), then we let $u$ be the vertex in $V$ that it corresponds to. For each pair $F,F'\in \fset$, with $F$ corresponding to $u_F$ and and $F'$ corresponding to $u_{F'}$, we define $\hat\delta(F,F')=\dist_{\ell}(u_F,u_{F'})$. 
As graph $G$ in the instance $(G,T,\delta)$ constructed in \Cref{subsec: instance} is an unweighted graph, it is easy to see that
\[\hat\delta(F,F')=\dist_{\ell}(u_F,u_{F'})\le \lcon(\phi(F),\phi(F'))+2.\]
As the mapping $\phi$ is random, $\hat{\delta}$ is also random, and so $\ex{\hat\delta(F,F')}\le O(1)\cdot  \delta(F,F') +2$.

From the properties of mapping $\phi$ in \Cref{lem:d-g}, we are guaranteed that such a solution $(\fset,\hat\delta)$ is a canonical solution. Moreover, from linearity of expectation, 
$$\ex{\vol(\fset,\hat\delta)}
=\ex{\sum_{(u,v)\in E}\hat\delta(F(u),F(v))}
=\sum_{(u,v)\in E}O\bigg(\delta(F(u),F(v))\bigg)+2
=O\bigg(\vol(\fset,\delta)\bigg)+O(n).
$$
Therefore, it follows that there exists a canonical solution $(\fset,\hat\delta)$, such that $\vol(\fset,\hat\delta)\le O(\vol(\fset,\delta)+n)$.
As we have shown in \Cref{subsec: canonical} that any canonical solution $(\fset,\hat\delta)$ with $|\fset|\le o(k\log^{1-\eps}k)$ satisfies that $\vol(\fset,\hat\delta)=\Omega(\eps n\log\log n)$, it follows that $\vol(\fset,\delta)=\Omega(\eps n\log\log n)$.
%As $G$ is an unweighted graph with $O(n)$ edges, the solution $(\fset,\hat\delta)$ incurs a stretch of $\Omega(\eps \log\log n)=\Omega(\eps \log\log k)$.
This implies that the integrality gap of (\textsf{LP-Metric}) is at least $\Omega(\eps \log\log n)$.
\qed






\subsection{Proof of \Cref{lem:d-g}}
\label{subsec: diam-girth}


In this subsection, we provide the proof of \Cref{lem:d-g}. We first consider the special case where $G$ is a tree, and then prove \Cref{lem:d-g} for the general case.
%Before we prove the statement in any graph, we first prove that in the case when $G$ is a tree, we can preserve the disance. 

\begin{lemma} \label{lem:embed-tree}
Let $(G,T,\ell)$ be an instance of $\zesn$ where $G$ is a tree. Let $(\fset,\delta)$ be an solution to it. Let $(\Vcon, \lcon)$ be the continuization of $G$. Then there exists a mapping $\phi: \fset \to \Vcon$, such that
\begin{itemize}
	\item for each terminal $t\in T$, if $F$ is the (unique) cluster in $\fset$ that contains $t$, then $\phi(F)=t$; and
	\item for every pair $F,F'\in \fset$,
	%we can randomly construct a projection $f$ from the nodes in $\delta$ to the veritces in $V$ such that: for any pair of nodes $F_1$ and $F_2$,
	$\lcon(\phi(F),\phi(F')) \le \delta(F,F')$.
\end{itemize}
%Given a 0-Extension instance $(G,T,\ell_e)$ where $G$ is a tree, and given a metric $\delta$ that contains all vertices in $T$, we can randomly construct a projection $f$ from the nodes in $\delta$ to the points (including vertices and points on the edge) such that: for any pair of nodes $F_1$ and $F_2$, $\ell(\phi(F_1),\phi(F_2)) \le \delta(F_1,F_2)$.
\end{lemma}

\begin{proof}
For each terminal $t\in T$, we denote by $F_t$ the cluster in $\fset$ that contains it, and set $\phi(F_t)=t$.
For each cluster $F\in \fset$ that does not contain any terminals, we define \[\nu(F)=\min\set{\frac{1}{2}\cdot \big(\delta(F,F_t)+\delta(F,F_{t'})-\delta(F_t,F_{t'})\big)\mid t,t'\in T}.\]
%As $\delta$ is a metric on $\fset$, $\nu(F)\ge 0$. 
Denote by $t_1,t_2$ the pair of terminals $(t,t')$ that minimizes $(\delta(F,F_t)+\delta(F,F_{t'})-\delta(F_t,F_{t'}))/2$.
As $G$ is a tree, there is a unique shortest path connecting $t_1$ to $t_2$ in $G$, and therefore there exists a unique point $p$ in $\Vcon$ (that lies on the $t_1$-$t_2$ shortest path in $\Vcon$) with $\lcon(p,t_1)=\delta(F,F_{t_1})-\nu(F)$ and $\lcon(p,t_2)=\delta(F,F_{t_2})-\nu(F)$. We set $\phi(F)=p$.
	
%For any node $s$ in $\delta$, let $\nu(s)$ be the maximum distance $\nu$ such that for any two terminals $t_1$ and $t_2$ in $T$, $\delta(s,t_1)+\delta(s,t_2) \ge \delta(t_1,t_2) + 2\nu$. Since $\delta$ is a metric, we have $\nu(s) \ge 0$. Moreover, since $\nu(S)$ is the maximum possible distance, there exist two terminals $t_1(s)$ and $t_2(s)$ such that $\delta(s,t_1(s))+\delta(s,t_2(s)) = \delta(t_1(s),t_2(s)) + 2\nu(s)$. Let $\phi(s)$ be the point in the tree path between $t_1(s)$ and $t_2(s)$ such that $\ell(\phi(s),t_1(s)) = \delta(s,t_1(s)) - \nu(s)$ and $\ell(\phi(s),t_2(s)) = \delta(s,t_2(s)) - \nu(s)$. 

%Before we prove the correctness, we first prove two proposition of the projection $f$.
 
\iffalse
\begin{claim} \label{clm:t-mid}
        Given three nodes $t,t_1,t_2$ on a tree such that $t$ is on the tree path between $t_1$ and $t_2$, for any node $t'$ on the tree, $t$ is either on the tree path between $t'$ and $t_1$ or the tree path between $t'$ and $t_2$.
    \end{claim}

    \begin{proof}
        Since the path between $t'$ and $t_1$ and the path between $t'$ and $t_2$ fully contain the path between $t_1$ and $t_2$. So $t$ must be on one of these two paths
    \end{proof}
\fi

\begin{claim} \label{clm:t-proj}
For every cluster $F\in \fset$ and every terminal $t\in T$, $\lcon(\phi(F),t) \le \delta(F,F_t) - \nu(F)$. 
\end{claim}

\begin{proof}
Note that the point $\phi(F)$ lies on the (unique) shortest path between a pair $t_1,t_2$ of terminals.
For $t\in \set{t_1,t_2}$, clearly the claim holds.
Consider any other terminal $t$. Clearly $\phi(F)$ lies on either the path connecting $t$ to $t_1$ or the path connecting $t$ to $t_2$. 
Assume without lose of generality that $\phi(F)$ is on the path connecting $t_1$ and $t$.
%By \Cref{clm:t-mid}, without lose of generality, suppose $\phi(F)$ is on the path between $t_1(F)$ and $t$. 
Since $G$ is a tree, $\lcon(\phi(F),t_1)+\lcon(\phi(F),t) = \lcon(t_1,t) = \dist_{\ell}(t_1,t)$. On the other hand, by definition of $\nu(F)$ and $\phi(F)$, $\delta(F,F_{t_1})+\delta(F,F_t) \ge \delta(F_{t_1},F_{t}) + 2 \cdot\nu(F)$ holds, and $\lcon(\phi(F),t_1) = \delta(F,F_{t_1}) - \nu(F)$. Therefore, $\delta(F,F_t) \ge \lcon(\phi(F),t) + \nu(F)$.
\end{proof}
    
We now show that, for every pair $F,F'\in \fset$, $\lcon(\phi(F),\phi(F')) \le \delta(F,F')$. 
Denote by $t_1,t_2$ the pair of terminals whose shortest path contains $\phi(F)$, and by $t'_1,t'_2$ the pair of terminals whose shortest path contains $\phi(F')$.
Assume without lose of generality that $\phi(F)$ is on the tree path between $\phi(F')$ and $t_1$, so $\lcon(\phi(F'),\phi(F)) + \lcon(\phi(F),{t_1}) = \lcon(\phi(F'),{t_1})$. On the other hand, from the definition of $\phi(F)$ and \Cref{clm:t-proj}, $\lcon(\phi(F),{t_1}) = \delta(F,F_{t_1}) - \nu(F)$ and $\lcon(\phi(F'),{t_1}) \le \delta(F',F_{t_1}) - \nu(F')$. Therefore, 
$$\lcon(\phi(F),\phi(F')) \le \delta(F',F_{t_1}) - \delta(F,F_{t_1}) - \nu(F') + \nu(F) \le \delta(F,F') + \nu(F) - \nu(F').$$ Similarly, $\lcon(\phi(F),\phi(F')) \le \delta(F,F') + \nu(F') - \nu(F)$. Altogether, $\lcon(\phi(F),\phi(F')) \le \delta(F,F')$.
\end{proof}


In the remainder of this subsection, we complete the proof of \Cref{lem:d-g}.
%Now we prove the general case. Let $D$ be the diameter and $\girth$ be the girth of the graph. 
Denote $g=\gir(G)$. 
Let $r$ be a real number chosen uniformly at random from the interval $[g/60,g/30]$, so $r\le g/30$ always holds.
%Let $r<\girth/30$ be a parameter whose value is to be set later. %Let $t^*$ be an arbitrarily chosen terminal in $T$. 
For each terminal $t\in T$, we denote by $F_t$ the cluster in $\fset$ that contains it, and set $\phi(F_t)=t$, so the first condition in \Cref{lem:d-g} is satisfied.


For each cluster $F\in \fset$, we define $A_F=\min\set{\delta(F,F_t)\mid t\in T}$.
%Given a point $s$, define $A_s$ as the distance between $s$ and its closest terminal. Let $r$ be a distance which is smaller than $\girth/30$ and we will decide it later, and let $t$ be an arbitrary terminal. 
We first determine the image $\phi(F)$ for all clusters $F$ with $A_F \le r$, in a similar way as \Cref{lem:embed-tree} as follows.
%\znote{to add some intuition about the change}

%For any $s$ such that $A_s \le r$, we use similar projection scheme as the tree case, but with some modification. We project $s$ as follows: 
Define \[\nu(F)=\min\set{\frac{1}{2}\cdot \bigg(\delta(F,F_t)+\delta(F,F_{t'})-\frac{\delta(F_t,F_{t'})}{2}\bigg)\mid t,t'\in T}.\]
%As $\delta$ is a metric on $\fset$, $\nu(F)\ge 0$. 
Denote by $t_1,t_2$ the pair $(t,t')$ that minimizes the above formula. We prove the following claim.
%let $\nu(F)$ be the maximum distance $\nu$ such that for any two terminals (can be the same) $t_1$ and $t_2$ in $T$, $\delta(F,t_1)+\delta(F,t_2) \ge \delta(t_1,t_2)/2 + 2\nu$. Let $t_1(S)$ and $t_2(S)$ be those two terminals that make the inequatlity equal.

\begin{claim} \label{clm:g-tree}
$\delta(F,F_{t_1}) + \delta(F,F_{t_2}) \le 4\cdot A_F \le 4r$.
\end{claim}
\begin{proof}
Let $t$ be the terminal such that $\delta(F,F_t) = A_F$. By definition of $\nu(F)$, 
$$\nu(F) \le \frac{\delta(F,F_t)+\delta(F,F_t)- \frac{1}{2}\cdot \delta(F_t,F_t)}{2} = A_F.$$ 
On the other hand, from triangle inequality, %since $\delta(F_{t_1},F_{t_2}) \le \delta(F,F_{t_1})+\delta(F,F_{t_2})$, 
$$\nu(F) = \frac{\delta(F,F_{t_1})+\delta(F,F_{t_2}) - \frac 1 2\cdot \delta(F_{t_1},F_{t_2})}{2} \ge \frac{\delta(F,F_{t_1})+\delta(F,F_{t_2})}{4}.$$ Altogether, $\delta(F,F_{t_1})+\delta(F,F_{t_2}) \le 4\cdot A_F$. 
\end{proof}

From \Cref{clm:g-tree}, $\delta(F_{t_1},F_{t_2}) \le \delta(F,F_{t_1}) + \delta(F,F_{t_2})\le 4r < \girth/3$. Therefore, there is a unique shortest path connecting $t_1$ to $t_2$ in $G$, as otherwise $G$ must contain a cycle of length at most $2\girth/3$, contradicting the fact that $\gir(G)=\girth$. %the following claim is close to \Cref{clm:t-proj}.
 
We then set $\phi(F)$ to be the point in $\Vcon$ that lies in the $t_1$-$t_2$ shortest path, such that $\lcon(\phi(F),t_1) = 2\cdot (\delta(F,F_{t_1}) - \nu(F))$ and $\lcon(\phi(F),t_2) = 2\cdot (\delta(F,F_{t_2}) - \nu(F))$. Note that, by definition of $t_1,t_2$,
\[
\lcon(\phi(F),t_1)+\lcon(\phi(F),t_2)=2\cdot \delta(F,F_{t_1}) +2\cdot \delta(F,F_{t_2}) - 4\cdot \nu(F)=\delta(F_{t_1},F_{t_2}).
\] 
%We first prove that when $A_F$ is  small, the projection is similar to the tree case.
We prove the following claim, that is similar to \Cref{clm:t-proj}.
    

    \begin{claim} \label{clm:g-proj}
        For every terminal $t$ with $\delta(F,F_t) \le 6r$, $\lcon(\phi(F),t) \le 2\cdot (\delta(F,F_t) - \nu(F))$.
    \end{claim} 

    \begin{proof}
From \Cref{clm:g-tree}, 
\[
\begin{split}
\delta(F_t,F_{t_1}) + \delta(F_t,F_{t_2}) + \delta(F_{t_1},F_{t_2}) & \le  \bigg(2\cdot \delta(F,F_t) + \delta(F,F_{t_1}) + \delta(F,F_{t_2})\bigg) + \delta(F_{t_1},F_{t_2}) \\ 
& \le  12 r + 4 r + 4 r + 4 r < \girth.
\end{split}
\] 
Therefore, the point $\phi(F)$ must lie on either the $t$-${t_1}$ shortest path or the $t$-${t_2}$ shortest path in $\Vcon$, as otherwise the union of $t$-${t_1}$ shortest path, $t$-${t_2}$ shortest path, and $t_1$-${t_2}$ shortest path contains a cycle of length less than $\girth$, a contradiction.  

Assume without loss of generality that $\phi(F)$ lies on the $t$-${t_1}$ shortest path, so 
$$\lcon(\phi(F),t) = \delta(F_t,F_{t_1}) - \lcon (\phi(F),{t_1}) = \delta(F_t,F_{t_1}) - 2(\delta(F,F_{t_1}) - \nu(F)).$$ By definition of $\nu(F)$, $\delta(F_t,F_{t_1}) \le 2\cdot (\delta(F,F_t)+\delta(F,F_{t_1}) - 2\cdot\nu(F))$. Therefore, $$\lcon(\phi(F),t) \le 2\bigg(\delta(F,t)+\delta(F,F_{t_1}) - 2\nu(F) - \big(\nu(F,F_{t_1}) -\nu(F)\big)\bigg) = 2(\delta(F,F_t)-\nu(F)).$$
    \end{proof}

%Like the tree case, we are ready to prove the correctness of the projection.

We now show in the next claim that the second condition in \Cref{lem:d-g} holds for pairs of clusters that are close in $\delta$.

\begin{claim} \label{clm:g-close}
For every pair $F,F'\in \fset$ with $A_{F}, A_{F'},\delta(F,F') \le r$,  $\lcon(\phi(F),\phi(F')) \le 2\cdot\delta(F,F')$. 
    \end{claim}

\begin{proof}
Since $\delta(F,F') < r$, from triangle inequality and \Cref{clm:g-tree}, 
$$\delta(F',F_{t_1})\le \delta(F',F)+\delta(F,F_{t_1})\le r+4r\le 5r.$$ 
Similarly, $\delta(F',F_{t_2})\le 5r$. 
Then from \Cref{clm:g-proj}, 
$$\lcon(\phi(F'),t_1) \le 2\cdot (\delta(F',F_{t_1}) - \nu(F')) \le 2\cdot \delta(F',F_{t_1}) \le  10r,$$ 
and symmetrically, $\lcon(\phi(F'),t_2) < 10r$. Therefore, $$\lcon(\phi(F'),t_1)+\lcon(\phi(F'),t_2)+\lcon(t_1,t_2) < 20r + 4r < \girth.$$
%So among $t_1(F),t_2(F)$ and $\phi(F')$, one vertex is on the shortest path between the other two vertices.
Therefore, the point $\phi(F)$ must lie either on the $\phi(F')$-$t_1$ shortest path or the $\phi(F')$-$t_2$ shortest path in $\Vcon$, as otherwise the union of $\phi(F')$-${t_1}$ shortest path, $\phi(F')$-${t_2}$ shortest path, and $t_1$-${t_2}$ shortest path in $\Vcon$ contains a cycle of length less than $\girth$, a contradiction. 

Assume without loss of generality that $\phi(F)$ lies on the $\phi(F')$-$t_1$ shortest path. Then
\[
\begin{split}
\lcon(\phi(F),\phi(F')) & =\lcon(\phi(F'),t_1(F))-\lcon(\phi(F),t_1(F))\\ 
& \le 2\bigg(\delta(F',t_1(F))-\nu(F')\bigg) - 2\bigg(\delta(F,t_1(F))-\nu(F)\bigg) \le 2\bigg(\delta(F,F')+\nu(F)-\nu(F')\bigg).
\end{split}
\]  
Similarly, $\lcon(\phi(F),\phi(F')) \le 2\big(\delta(F,F')+\nu(F')-\nu(F)\big)$. Altogether, $\lcon(\phi(F),\phi(F')) \le 2\delta(F,F')$.
\end{proof}


We now complete the construction of the mapping $\phi$, by specifying the images of all clusters $F\in \fset$ with $A_F> r$. Let $t^*$ be an arbitrarily chosen terminal in $T$. For all clusters $F\in \fset$ with $A_F > r$, we simply set $\phi(F) = t^*$. 
%Now we give the complete projection scheme. We first choose an arbitrary terminal $t$ and randomly choose the distance $r$ between $\girth/60$ and $\girth/30$.
%Let $r$ be a real number chosen uniformly at random from the interval $[g/60,g/30]$. We have constructed the image $\phi(F)$ for all clusters $F$ with $A_F \le r$. 
%we use the projection scheme we described before.


It remains to show that the second condition in \Cref{lem:d-g} holds for all pairs $F,F'\in \fset$. Consider a pair $F,F'$.
%For any two points $F$ and $F'$, 
Assume first that $\delta(F,F')>\girth/60$. Then $$\lcon(\phi(F),\phi(F')) \le \dia(G) \le \frac{60 \cdot \dia(G)}{\gir(G)}\cdot \delta(F,F').$$ 
Assume now that $\delta(F,F') \le \girth/60$, and without loss of generality that $A_{F} \le A_{F'}$. 
Note that, from triangle inequality, for every terminal $t\in T$, $\delta(F',F_t)\le \delta(F,F_t)-\delta(F,F')$. This implies that $A_{F'}-A_{F} \le \delta(F,F')$. Therefore, the probability that the random number $r$ takes value from the interval $[A_{F}, A_{F'}]$ is at most $\frac{\delta(F,F')}{\girth/60}$. %In this case, $\lcon(\phi(F),\phi(F')) \le \dia(G)$. 
Note that, if $r\le A_{F}$, then $\phi(F)=\phi(F')=t^*$ and $\lcon(\phi(F),\phi(F'))=0$.
%If they are both larger than $r$, then they are both projected to $t$ and thus $\lcon(\phi(F),\phi(F'))=0$. 
And  if $r\ge A_{F'}$, then from \Cref{clm:g-close}, $\lcon(\phi(F),\phi(F')) \le 2\delta(F,F')$. Altogether,
$$\ex{\lcon(\phi(F),\phi(F'))} \le \frac{60\cdot \dia(G)}{\gir(G)} \cdot\delta(F,F') + 2 \delta(F,F') \le O\bigg(\frac{\dia(G)}{\gir(G)}\bigg)\cdot \delta(F,F').$$

%    Finally, notice that in the current scheme, we are allow to project the point to not only the vertices in the graph, but also the points between two vertices. To avoid this, we can project $s$ to the closest vertices to $\phi(F)$, we call it $f'(F)$. It clear that for any two points $F$ and $F'$, $\ell(f'(F),f'(F')) \le \ell(\phi(F_1),\phi(F_2)) + 1$.
